\notebook{11-binary-tricks}
The questions of this section concern the integers of Appendix~\ref{notes:chap.integersInBinary}
and the books of Section~\ref{notes:sec.fasterBook} that hold the occupied prices of a side as one of them.
\codehl{TickArrayBook} and \codehl{BitmapBook} are classes of \codehl{unito26.lob.orderbook},
subclasses of \codehl{AggregateBook} that inherit its matching.
\codehl{TickArrayBook(floor, width)} indexes the band of $\bandWidth$ prices
from its floor $\bandFloor$ to its \codehl{ceiling}, $\bandFloor + \bandWidth - 1$;
it keeps the size at $p$ on each side in the entry $p - \bandFloor$ of a list,
and the occupied prices of the side \codehl{d} as the set bits of the integer \codehl{\_bits[d]}.
\codehl{BitmapBook(floor)} keeps its sizes in the dictionaries of \codehl{AggregateBook},
and its occupied prices in \codehl{\_bits[d]} as well.
\codehl{book.best\_bid\_price} and \codehl{book.best\_ask\_price} are
\codehl{book.best\_price(BUY)} and \codehl{book.best\_price(SELL)},
and \codehl{book.spread} is the second less the first, or \codehl{None} when either is \codehl{None}.
\codehl{ReportedDepth} is an integer type of \codehl{unito26.lob.messages},
and the functions of \codehl{unito26.lob.binary\_gaps} are imported by name.

\begin{question}
\label{question.tickArrayIntegersAfterACrossingBuy}
A limit buy that crosses the spread reaches a \codehl{TickArrayBook} of ten ticks.
The book inherits the matching of \codehl{AggregateBook}, quoted below;
\codehl{resting\_price} returns the limit price of a limit order,
and for a market order the price of its last fill.
\begin{lstlisting}[caption={\protect\codehl{unito26.lob.orderbook.AggregateBook.submit}, abridged.}, label=listing.tickArrayIntegersSubmit]
def submit(self, message: Message, record: bool) -> SubmitResult | None:
# ...
    remaining = message.size
    direction = message.direction
    limit_price = message.price
# ...
    while remaining > 0:
        best = self.best_price(-direction)
# ...
        if best is None or best * direction > limit_price * direction:
            break
        resting = self.size_at(-direction, best)
        traded = min(remaining, resting)
        remaining -= traded
        self.set_size(-direction, best, resting - traded)
# ...
    rest_price = self.resting_price(limit_price, last_fill_price)
    if remaining > 0 and rest_price is not None:
        resting = self.size_at(direction, rest_price) + remaining
        self.set_size(direction, rest_price, resting)
\end{lstlisting}
\codehl{TickArrayBook} writes a size, and finds the best price, as follows.
\begin{lstlisting}[caption={\protect\codehl{unito26.lob.orderbook.TickArrayBook}: a write, and the best price, abridged.}, label=listing.tickArrayIntegersWrite]
def set_size(self, direction: int, price: int, size: int) -> None:
    slot = price - self.floor
# ...
    self._sizes[direction][slot] = size
    bit = 1 << (price - self.floor if direction == BUY else self.ceiling - price)
    if size > 0:
        self._bits[direction] |= bit
    else:
        self._bits[direction] &= ~bit

def best_price(self, direction: int) -> int | None:
    bits = self._bits[direction]
    if not bits:
        return None
    index = bits.bit_length() - 1
    return self.floor + index if direction == BUY else self.ceiling - index
\end{lstlisting}
\begin{lstlisting}[caption={The two integers after a crossing buy}, label=listing.tickArrayIntegersRun]
book = TickArrayBook(floor=200, width=10)
for price, size in ((203, 50), (202, 40), (200, 70)):
    book.set_size(BUY, price, size)
for price, size in ((205, 60), (206, 30), (208, 90)):
    book.set_size(SELL, price, size)
book.apply(limit_order(1.0, 100, 206, BUY), record=False)
print(book._bits[BUY], book._bits[SELL],
      book.best_bid_price, book.best_ask_price)
\end{lstlisting}
What does the snippet print?
\begin{choices}
 \choice \codehl{77 256 206 208}
 \choice \codehl{77 4 206 208}
 \correctchoice \codehl{77 2 206 208}
 \choice \codehl{13 2 203 208}
 \choice \codehl{77 26 206 205}
\end{choices}
\explanation{The ceiling is 209, so the asks at 205, 206 and 208 are bits 4, 3 and 1, the integer 26.
By Proposition~\ref{notes:prop.decompositionOfLimitOrder}
the buy takes 60 at 205 and 30 at 206, clearing bits 4 and 3,
and rests its last 10 at 206, bit 6 of the bids: $13 + 64 = 77$.
Counted up from the floor, as the bids are, the ask left at 208 would be bit 8, the integer 256.}
\end{question}

\begin{question}
\label{question.bitmapBestAskFromTheWrongEnd}
\codehl{BitmapBook} writes a size as below,
the bit of a price counted up from the floor on both sides.
\begin{lstlisting}[caption={\protect\codehl{unito26.lob.orderbook.BitmapBook.set\_size}, abridged.}, label=listing.bitmapBestAskSetSize]
def set_size(self, direction: int, price: int, size: int) -> None:
# ...
    super().set_size(direction, price, size)
    bit = 1 << (price - self.floor)
    if size > 0:
        self._bits[direction] |= bit
    else:
        self._bits[direction] &= ~bit
\end{lstlisting}
Two variants of its \codehl{best\_price} read the best price off the integer,
and a market buy for the whole ask side arrives,
priced at \codehl{sys.maxsize}, above every price of the band,
matched as in Listing~\ref{listing.tickArrayIntegersSubmit}.
\begin{lstlisting}[caption={Alternative A}, label=listing.bitmapBestAskA]
class Book(BitmapBook):
    def best_price(self, direction):
        return self.floor + self._bits[direction].bit_length() - 1
\end{lstlisting}
\begin{lstlisting}[caption={Alternative B}, label=listing.bitmapBestAskB]
class Book(BitmapBook):
    def best_price(self, direction):
        bits = self._bits[direction]
        if direction == BUY:
            return self.floor + bits.bit_length() - 1
        return self.floor + (bits & -bits).bit_length() - 1
\end{lstlisting}
\begin{lstlisting}[caption={Run once with each alternative}, label=listing.bitmapBestAskRun]
book = Book(floor=95)
for price, size in ((98, 40), (97, 60)):
    book.set_size(BUY, price, size)
for price, size in ((101, 30), (102, 50), (105, 20)):
    book.set_size(SELL, price, size)
before = book.best_ask_price
book.apply(market_order(1.0, 100, BUY), record=False)
print(before, book.best_ask_price, book.spread)
\end{lstlisting}
What happens?
\begin{choices}
 \choice Both print \codehl{101 None None}.
 \correctchoice A prints \codehl{105 94 -4}; B prints \codehl{101 94 -4}.
 \choice A prints \codehl{105 None None}; B prints \codehl{101 None None}.
 \choice Both print \codehl{101 94 -4}.
 \choice A prints \codehl{105 95 -3}; B prints \codehl{101 95 -3}.
\end{choices}
\explanation{Counted up from the floor, the best ask is the lowest set bit,
which B finds by Proposition~\ref{notes:prop.lowestSetBit};
A returns the highest ask, 105, and its market buy fills from the top of the side.
On the emptied side \codehl{0 \& -0} is $0$, whose \codehl{bit\_length} is $0$,
so both return $95 - 1$, a price below the bid, and no error.}
\end{question}

\begin{question}
\label{question.occupancyWrittenByAddition}
Two variants of the write of Listing~\ref{listing.tickArrayIntegersWrite}
update the integer of occupied prices with other operators.
After a message,
\codehl{book.best\_bid\_size} is the entry of the list of bid sizes at the best bid.
\begin{lstlisting}[caption={Alternative A}, label=listing.occupancyByAdditionA]
class Book(TickArrayBook):
    def set_size(self, direction, price, size):
        self._sizes[direction][price - self.floor] = size
        bit = 1 << (price - self.floor if direction == BUY
                    else self.ceiling - price)
        if size > 0:
            self._bits[direction] += bit
        else:
            self._bits[direction] -= bit
\end{lstlisting}
\begin{lstlisting}[caption={Alternative B}, label=listing.occupancyByAdditionB]
class Book(TickArrayBook):
    def set_size(self, direction, price, size):
        sizes = self._sizes[direction]
        slot = price - self.floor
        bit = 1 << (slot if direction == BUY else self.ceiling - price)
        if (sizes[slot] > 0) != (size > 0):
            self._bits[direction] ^= bit
        sizes[slot] = size
\end{lstlisting}
\begin{lstlisting}[caption={Run once with each alternative}, label=listing.occupancyByAdditionRun]
book = Book(floor=95, width=12)
for price, size in ((99, 100), (98, 60)):
    book.set_size(BUY, price, size)
book.apply(market_order(1.0, 40, SELL), record=False)
print(book.best_bid_price, book.best_bid_size)
\end{lstlisting}
What happens?
\begin{choices}
 \choice Both print \codehl{99 60}.
 \choice A prints \codehl{99 60}; B prints \codehl{98 60}.
 \choice A prints \codehl{100 60}; B prints \codehl{99 60}.
 \correctchoice A prints \codehl{100 0}; B prints \codehl{99 60}.
 \choice A prints \codehl{100 0}; B prints \codehl{98 60}.
\end{choices}
\explanation{Addition agrees with $\bitOr$ only on a bit that is clear.
The market sell leaves 60 at 99, a write to a price whose bit 4 is set,
and $24 + 16 = 40$ carries into bit 5:
the best bid becomes 100, where nothing rests.
B flips the bit only when the price passes between empty and occupied,
which a partial fill does not.}
\end{question}

\begin{question}
\label{question.occupancyClearedWithMinusForTilde}
In the variant below of the write of Listing~\ref{listing.tickArrayIntegersWrite},
the line that clears the bit of an emptied price has a minus where the package has a tilde.
\codehl{withdrawal(t, q, p, d)} removes $q$ shares from the size at $p$ on the side $d$
through \codehl{set\_size},
and \codehl{book.occupied\_levels(d, L)} walks the set bits of the side $d$ from the highest down,
returning for each the pair of its price and the size the list holds there,
at most $L$ of them.
\begin{lstlisting}[caption={A minus for a tilde}, label=listing.occupancyClearedWithMinus]
class Book(TickArrayBook):
    def set_size(self, direction, price, size):
        self._sizes[direction][price - self.floor] = size
        bit = 1 << (price - self.floor if direction == BUY
                    else self.ceiling - price)
        if size > 0:
            self._bits[direction] |= bit
        else:
            self._bits[direction] &= -bit

book = Book(floor=95, width=12)
for price, size in ((99, 100), (98, 60), (96, 40)):
    book.set_size(BUY, price, size)
book.apply(withdrawal(1.0, 60, 98, BUY), record=False)
print(book.best_bid_price, book.occupied_levels(BUY, ReportedDepth(3)))
\end{lstlisting}
What does the snippet print?
\begin{choices}
 \choice \codehl{98 [(98, 0)]}
 \choice \codehl{99 [(99, 100), (98, 0), (96, 40)]}
 \choice \codehl{99 [(99, 100), (96, 40)]}
 \choice \codehl{99 [(99, 100)]}
 \correctchoice \codehl{99 [(99, 100), (98, 0)]}
\end{choices}
\explanation{By Definition~\ref{notes:def.twosComplementSequence},
$-8 = \neg 7$ has every bit from 3 upwards set and the three below clear,
so $26 \bitAnd (-8) = 24$:
the emptied price at bit 3 stays, and every deeper one, 96 among them, is cleared.
The tilde gives $\neg 8 = -9$, which clears bit 3 alone,
Proposition~\ref{notes:prop.twosComplement}.}
\end{question}

\begin{question}
\label{question.gapsOfTheUncutBidOccupancy}
A gap of a side is a maximal run of empty prices lying between two occupied ones.
Two functions of the package, and a one-liner, are applied to the bid integer of a
\codehl{TickArrayBook}.
\codehl{bin(n)} is the string \codehl{"0b"} followed by the binary digits of $n > 0$, the highest first,
and \codehl{s.count(t)} is the number of non-overlapping occurrences of \codehl{t} in \codehl{s}.
\begin{lstlisting}[caption={\protect\codehl{unito26.lob.binary\_gaps}: two functions, their docstrings elided.}, label=listing.gapsUncutFunctions]
def measure_largest_binary_gap(n: int) -> int:
# ...
    count = 0
    while (n & (n + 1)) != 0:
        count += 1
        n |= n >> 1
    return count


def count_binary_gaps(n: int) -> int:
# ...
    if not n:
        return 0
    return (n & ~(n << 1)).bit_count() - 1
\end{lstlisting}
\begin{lstlisting}[caption={Three readings of the bid integer}, label=listing.gapsUncutRun]
book = TickArrayBook(floor=95, width=12)
for price, size in ((104, 10), (103, 20), (100, 30), (99, 40)):
    book.set_size(BUY, price, size)
bits = book._bits[BUY]
print(count_binary_gaps(bits), bin(bits)[2:].count("10"),
      measure_largest_binary_gap(bits))
\end{lstlisting}
What does the snippet print?
\begin{choices}
 \correctchoice \codehl{1 2 4}
 \choice \codehl{1 1 2}
 \choice \codehl{1 2 2}
 \choice \codehl{1 1 4}
 \choice \codehl{2 2 4}
\end{choices}
\explanation{The integer is $1100110000$ in binary:
one gap of two ticks, 101 and 102, above four trailing zeros, the empty prices from 95 to 98.
\codehl{count\_binary\_gaps} keeps one bit per run of ones, Proposition~\ref{notes:prop.runsOfOnes},
and trailing zeros make no run;
the one-liner also counts the step from the lowest one into them,
and the flood fills them in four passes.
\codehl{span\_bits} shifts the window down to an odd integer first.}
\end{question}

\begin{question}
\label{question.askWindowsOfTheTwoBitmaps}
The gap statistics of a side are read on a window of its integer,
the bits from the touch to the deepest reported level, shifted down so that the lower end of the window is bit 0.
For the ask side and the reported depth $L$, \codehl{span\_bits(SELL, L)} returns
\codehl{keep\_lowest\_set\_bits(bits, L)} in \codehl{BitmapBook}
and \codehl{keep\_highest\_set\_bits(bits, L)} in \codehl{TickArrayBook},
\codehl{bits} being the ask integer of each.
\codehl{count\_trailing\_zeros(n)} is the number of zeros below the lowest set bit of $n > 0$.
\begin{lstlisting}[caption={\protect\codehl{unito26.lob.binary\_gaps}: three functions, their docstrings elided.}, label=listing.askWindowsFunctions]
def discard_trailing_zeros(n: int) -> int:
# ...
    return n >> count_trailing_zeros(n) if n else 0


def keep_highest_set_bits(n: int, count: int) -> int:
# ...
    if n.bit_count() <= count:
        return discard_trailing_zeros(n)
    low, high = 0, n.bit_length() - 1  # (n >> low).bit_count() >= count throughout
    while low < high:
        mid = (low + high + 1) // 2
        if (n >> mid).bit_count() >= count:
            low = mid
        else:
            high = mid - 1
    return n >> low


def keep_lowest_set_bits(n: int, count: int) -> int:
# ...
    rest = n
    for _ in range(count):
        rest &= rest - 1
    return discard_trailing_zeros(n ^ rest)
\end{lstlisting}
\begin{lstlisting}[caption={The window of three ask levels, in either book}, label=listing.askWindowsRun]
asks = {101: 30, 102: 50, 105: 20, 106: 40}
bitmap = BitmapBook(floor=95)
array = TickArrayBook(floor=95, width=12)
for price, size in asks.items():
    bitmap.set_size(SELL, price, size)
    array.set_size(SELL, price, size)
depth = ReportedDepth(3)
print(bin(bitmap.span_bits(SELL, depth)),
      bin(array.span_bits(SELL, depth)))
\end{lstlisting}
What does the snippet print?
\begin{choices}
 \choice \codehl{0b11001 0b11001}
 \choice \codehl{0b10011 0b10011}
 \choice \codehl{0b110011 0b110011}
 \correctchoice \codehl{0b10011 0b11001}
 \choice \codehl{0b11001 0b10011}
\end{choices}
\explanation{\codehl{BitmapBook} counts the asks up from the floor, bits 6, 7, 10 and 11,
and three passes of \codehl{rest \&= rest - 1} clear the lowest three, which \codehl{n \^{} rest} recovers:
101 lands on bit 0.
\codehl{TickArrayBook} counts them down from the ceiling, bits 5, 4, 1 and 0,
and the bisection of Proposition~\ref{notes:prop.rankByShift} cuts at bit 1:
105 lands on bit 0, and each window is the other read backwards.}
\end{question}

\begin{question}
\label{question.lowestSetBitCostsTheBand}
\codehl{BitmapBook} finds its best price as below,
and \codehl{TickArrayBook} as in Listing~\ref{listing.tickArrayIntegersWrite}.
CPython stores a non-negative integer as its digits in base $2^{30}$, together with their number.
\begin{lstlisting}[caption={\protect\codehl{unito26.lob.orderbook.BitmapBook.best\_price}}, label=listing.lowestSetBitCostsBitmap]
def best_price(self, direction: int) -> int | None:
    bits = self._bits[direction]
    if not bits:
        return None
    if direction == BUY:
        return self.floor + bits.bit_length() - 1
    return self.floor + (bits & -bits).bit_length() - 1
\end{lstlisting}
For a width $\bandWidth$, \codehl{width} below,
two books hold the same two asks, one a tick above the floor and one $\bandWidth - 1$ ticks above it.
\begin{lstlisting}[caption={The two books of a width}, label=listing.lowestSetBitCostsBooks]
bitmap = BitmapBook(floor=0)
array = TickArrayBook(floor=0, width=width)
for price in (1, width - 1):
    bitmap.set_size(SELL, price, 100)
    array.set_size(SELL, price, 100)
\end{lstlisting}
\codehl{bitmap.best\_ask\_price} and \codehl{array.best\_ask\_price} are read many times,
for $\bandWidth = 10^{3}, 10^{4}, 10^{5}, 10^{6}$,
and the time of one read is recorded.
Which statement holds?
\begin{choices}
 \correctchoice A read of \codehl{array} takes about the same time at every width;
  a read of \codehl{bitmap} grows with $\bandWidth$, about in proportion to it at the larger widths,
  and at $\bandWidth = 10^{6}$ is many times slower.
 \choice Both reads take about the same time at every width:
  each is a fixed number of operations on one integer.
 \choice Both reads grow about in proportion to $\bandWidth$:
  the two integers are equally wide.
 \choice A read of \codehl{bitmap} takes about the same time at every width,
  and a read of \codehl{array} grows with $\bandWidth$:
  its best ask is the highest set bit, at the far end of an integer of $\bandWidth$ bits,
  where that of \codehl{bitmap} is bit 1.
\end{choices}
\explanation{\codehl{bit\_length} reads the stored number of digits and the highest one, in $O(1)$;
the negation and $\bitAnd$ build new integers digit by digit, in $O(\bandWidth)$,
wherever the lowest set bit lies,
Section~\ref{notes:sec.binaryCosts}.
Counting the asks down from the ceiling makes the best ask the highest set bit.}
\end{question}
